\paragraph{Changing the path and its constants}
The path datum is a function of the tuple of residue matrices on its
edges. A change of basis at a node acts by the determinant characters
of the edges. The datum vanishes unless the vertical columns are
split injections. Fix the matrices outside a turn and transfer the
data of its two edges by the inverse rows and the first columns of
an invertible completion. Changing the completion changes the basis
of the column frame on the right. Its determinant character cancels
\(\chi(\det g)\) in the matching relation. Thus the transferred weight
is independent of the completion and obeys the same transformation law.
On the diagonal, a seesaw moves the period to
\(Y^{j+1}_j=\widetilde Y^{j+1}_j\). Taking the adjoint of the kernel
for the extra line applies a finite Fourier transform to its column.
The two edges on each path now fit \eqref{eq:exchange-square}.
After exchange, a second seesaw restores the form of the period.
The number of summands in a coefficient grows by at most \(N_\mathrm{col}\),
the number of possible residue columns. A completion contributes
a root of unity, and the line Fourier transform only multiplies by
additive-character roots of unity before summing over one column.
All \(m^2\) turns therefore cost at most \(N_\mathrm{col}^{m^2}\) in
term counts. The product of the nonzero denominator factors \(c_{ij}\)
from these turns is \(c_{\mathrm{path}}\).
\leanlinks{\leansrc{New/Determinant/AuxiliaryForms/Turns}{The node transformation law after changing the completion}
  \leanlinksep \leansrc{New/Determinant/AuxiliaryForms/DiagonalTurns}{The line on the diagonal and the composition of the turns}}

\paragraph{Denominators and sizes on the final path}
Put \(n_m=\dim X_m\) and \(d_b=b r^4\). Every diagram on the final path
has at most \(d_b\) boxes, and every rank is at most \(2n_m\).
The finite datum of an edge is the indicator of one residue matrix.
To return it to standard data by finite Fourier expansion, clear
the numbers of matrix residues at the finite level \(\pi_M^b\) and
at the conjugate places. Factors from the twist places are units
at \(\ell\nmid N_\xi\). A further \(M^b d\) per box clears the
derivatives of the integral vacuum. It follows that the common
integer
\[
 D_\mathrm{edge}=(M^b)^{2f(2n_m)^2}(M^b d)^{d_b}
\]
makes the operator on each edge preserve integrality.
At the endpoint the ranks differ by one and the highest weights
interlace. Expanding the highest weight vectors in branching minors
gives an integer branching matrix. Compose this with the Weyl
duality pairing to obtain the matrix for the endpoint pairing.
The basis of the last column at one endpoint is multiplied by
the \(m\)-th power of
\(u=\overline{t_{\mathfrak m}\alpha_\mathrm{W}}/(t_{\mathfrak m}\alpha_\mathrm{W})\),
where \(t_{\mathfrak m}\) generates \(\mathfrak m=\pi_M^bQ\) and
\(\alpha_\mathrm{W}\) is the vacuum idele. Since \(t_{\mathfrak m}\)
divides \(M^b N_\xi\), and \(q_1\alpha_\mathrm{W}\) and
\(\alpha_\mathrm{W}^{-1}\) are integral, \(M^b q_1\) clears both \(u\)
and \(u^{-1}\) at \(\ell\nmid N_\xi\). If \(e_\tau^\mathrm{last}\)
is the absolute determinant exponent, this costs
\(\prod_\tau(M^b q_1)^{m e_\tau^\mathrm{last}}\).
Include the orders of the class stabilizers in one denominator
\(D_\mathrm{end}\). Adding logarithms gives \(\log D_\mathrm{end}\le x^{98}b^2\).
There are \(m\) edges on each side, so the common integer
\[
 D_\mathrm{path}=D_\mathrm{end}D_\mathrm{edge}^{2m}
       \le \exp(x^{99}b^2)
\]
makes the final path sum integral. The path coefficients are sums
of roots of unity and add no denominator.
\leanlinks{\leansrc{New/Determinant/AuxiliaryForms/LastPath}{The integer branching matrix for interlacing weights}
  \leanlinksep \leansrc{New/Determinant/AuxiliaryForms/SpanningForms}{Clearing the final path denominators}}

For the trivial twist put \(G=\exp(x^{99}b^2)\).
Substituting \(d_b\) boxes and rank at most \(2n_m\) in the jet norm
bound of Section~\ref{sec:upper-bound} bounds each edge's conjugate
norm multiplier by \(G\). If \(B_-,B_+\) bound the initial squared
norms, the endpoint norms are at most \(G^m\sqrt{B_-}\) and
\(G^m\sqrt{B_+}\). The branching matrix has operator bound
\((d_b^{d_b})^3\). There are \(2m+1\) nodes whose class counts are
each at most \(G\), and at most \(N_\mathrm{col}^{m^2}N_\mathrm{mat}^{2m}\)
root-of-unity terms, where \(N_\mathrm{mat}\) bounds the number of
residue matrices per edge. Summing the bounds gives
\[
 |\sigma(\Xi)|\le
 N_\mathrm{col}^{m^2}N_\mathrm{mat}^{2m}(d_b^{d_b})^3G^{4m+1}\sqrt{B_-B_+}
 \le\exp(x^{200}b^2)\sqrt{B_-B_+}.
\]
The last inequality uses the admissible bounds for \(m,f,r,\log M\).
When testing with a dual delta form, \(\sqrt{B_-}\) is the norm of
its value. At infinity \(D_\mathrm{path}\) and the dual hook number
together cost at most \(\exp(x^{201}b^2-x^{200}b^2)\).
Division by the dual norm bounds the resulting vector norm by
\(\exp(x^{201}b^2)\) times the initial norm.
At deep level the class stabilizers are trivial and delta forms
span. Use short lattice bases in the Weyl representation and
polynomial evaluations at roots of unity to span the values on
each class. These give integral generators whose conjugate
squared norms are at most \(G\). Apply the surjection \(s\) to
obtain both bounds for the auxiliary forms.
\leanlinks{\leansrc{New/Determinant/AuxiliaryForms/SpanningForms}{Class counts, term counts, and the dual estimate}
  \leanlinksep \leansrc{New/UpperBound/Size/ShortColumns}{Spanning by short delta forms}}
